\noindent\small Each cell is a count over the records (or studies) behind the finding. The per-class validity tests of each row are Fig.~\ref{fig:findings} and Table~\ref{tab:vcov}, not a column here: the 11 classes are never summed, because each licenses a different, strictly stronger reading than the base one and none is a prerequisite for it, and eleven fractions do not fit a table cell. \emph{Unused data}: studies evaluating on a cohort or dataset unused for fitting or tuning, not necessarily an external clinical site. \emph{Split stated}: a patient-, slide- or subject-level split. \emph{Uncertainty}: repeated runs, seeds, confidence intervals or significance tests (resampling and seed replication are not distinguished in the coding). These three and \emph{Preprint} are study-level fields of the quality table (supplement); an em-dash marks a property that does not apply to the records of that finding, and unk the studies for which it could not be determined. \emph{Model families}: distinct model lineages (CLIP-style dual encoders, DINO-style self-supervised encoders, multimodal LLMs, task-specific classifiers, or unassigned) among the records, a coarse proxy for dependence between studies; \emph{distinct first authors}, the proxy the corpus-breadth tier uses, is a second such proxy. \emph{Rec.\ excl.\ preprints}: records from peer-reviewed studies only, with the corpus-breadth tier the row receives on those records alone (n/a when none remains). \emph{Abstract-only studies}: coded from the abstract because the full text could not be retrieved. \emph{Defect-touched studies}: studies with a recorded critical defect at least one of whose records is flagged, a record being flagged when the defect touches its own claim or its applicability is uncertain, that is, a defect that undermines the representation-level result (probe or dictionary fitted and evaluated on the same data, test-set training, leakage across split units, circular dependence estimates, selection on the reported metric, a single-case fitting block; the judgement is recorded per record in the released coding sheets). Every other column of this table, and every count in \S\ref{sec:findings} and Fig.~\ref{fig:findings}, is computed on all records, flagged or not; the parenthesis gives the corpus-breadth tier the row would receive with the flagged records set aside (the screen as a sensitivity, supplement S6). \emph{Dir.}: direction, read relative to the proposition of the row, so that for F4 and F8a a null result counts as positive. The leading triple is over studies, a study counting as positive only when every record it contributes to the row is positive and as mixed when its records disagree; the record triple follows it. A record can bundle several outcomes, so the study is the more interpretable unit. \emph{Unflagged} abbreviates \emph{not flagged for one of the five result-invalidating defect classes of Table~\ref{tab:defects}}; it is a screen against a short enumerated list and not a judgement that a study is at low risk of bias. The profile is a structured appraisal, not a formal risk-of-bias instrument, and the corpus-breadth tier is a descriptive statement about the corpus rather than a judgement of certainty.
